\subsection*{CU-25: Finalizar la partida y aplicar la progresión}
\addcontentsline{toc}{subsection}{CU-25: Finalizar la partida y aplicar la progresión}
\label{cu-25}

\begin{fichacu}{CU-25}{Finalizar la partida y aplicar la progresión}
  \crudFila{Responsable}{Ángel Gabriel Aguilar Hernández}
  \crudFila{Actor principal}{Servidor de partidas}
  \crudFila{Actores de sistema}{Base de datos}
  \crudFila{Actores interesados}{Jugador}
  \crudFila{Prototipos}{2g, 2h, 11a, 11b, 15c, 15d, 17d, 12g}
  \crudFila{Objetivo}{Cerrar la partida, determinar el resultado de cada participante y aplicar todo lo que la partida cambia en la cuenta: elo, división, experiencia, nivel, monedas y contadores.}
  \crudFila{Disparador}{Un peón llega a su fila de meta, un reloj llega a cero, un jugador se rinde o abandona, o los jugadores acuerdan tablas.}
  \textbf{Precondiciones} & \cuCelda{\begin{cureglas}
      \item \textbf{\mbox{PRE-01}} Existe una \textit{Partida} con estado \texttt{EN\_\allowbreak{}CURSO}.
      \item \textbf{\mbox{PRE-02}} Se ha producido una condición de término conforme a \mbox{RN-01}.
      \item \textbf{\mbox{PRE-03}} El Servidor de partidas tiene conexión disponible con la Base de datos.
    \end{cureglas}} \\ \hline
  \textbf{Flujo normal} & \cuCelda{\begin{cuenum}[start=1]
      \item El Servidor de partidas detecta la condición de término y determina la forma de término: \texttt{META}, \texttt{TIEMPO\_\allowbreak{}AGOTADO}, \texttt{RENDICION}, \texttt{ABANDONO} o \texttt{TABLAS} (\mbox{RN-01}).
      \item El Servidor de partidas determina el resultado de cada \textit{Participacion}: ganada, perdida o tablas, y la posición en el modo de cuatro jugadores (\mbox{RN-02}).
      \item El Servidor de partidas comprueba si la partida es clasificatoria, es decir, si nació del \mbox{CU-17} y no de una sala privada ni de la IA (\mbox{RN-03}). \textit{(Ver \mbox{FA-01})}
      \item El Servidor de partidas inicia una transacción sobre la Base de datos.
      \item El Servidor de partidas actualiza la \textit{Partida}: estado \texttt{FINALIZADA}, \texttt{fecha\_\allowbreak{}fin}, forma de término y el ganador, si lo hay. \textit{(Ver \mbox{EX-01})}
      \item El Servidor de partidas escribe en cada \textit{Participacion} su resultado, su posición y el tiempo de reloj que le quedaba.
      \item El Servidor de partidas calcula el elo nuevo de cada participante a partir del elo con el que entró, guardado en su \textit{Participacion}, y del resultado (\mbox{RN-04}).
      \item El Servidor de partidas escribe en cada \textit{Participacion} el elo resultante y la diferencia, para que la pantalla de fin y el historial puedan mostrarla sin recalcularla (\mbox{RN-05}).
    \end{cuenum}} \\
   & \cuCelda{\begin{cuenum}[start=9]
      \item El Servidor de partidas actualiza la \textit{EstadisticaModo} de cada participante en el \textit{Modo} de la partida: incrementa \texttt{partidas\_\allowbreak{}jugadas}, incrementa \texttt{partidas\_\allowbreak{}ganadas} o \texttt{partidas\_\allowbreak{}perdidas}, escribe \texttt{puntos\_\allowbreak{}elo}, actualiza \texttt{elo\_\allowbreak{}maximo} si procede, recalcula \texttt{racha\_\allowbreak{}actual} y \texttt{mejor\_\allowbreak{}racha}, suma el \texttt{tiempo\_\allowbreak{}total\_\allowbreak{}jugado} y suma las \texttt{barreras\_\allowbreak{}colocadas} contando las \textit{Jugada} de tipo \texttt{MURO} de ese jugador (\mbox{RN-06}). \textit{(Ver \mbox{EX-01})}
      \item El Servidor de partidas comprueba si el elo nuevo cruza un umbral de división, hacia arriba o hacia abajo, respetando las partidas de margen (\mbox{RN-07}). \textit{(Ver \mbox{FA-02})}
      \item El Servidor de partidas calcula la experiencia obtenida por cada participante y la suma a su \textit{Usuario}. \textit{(Ver \mbox{FA-03})}
      \item El Servidor de partidas calcula las monedas obtenidas y crea un \textit{MovimientoMoneda} de tipo \texttt{PARTIDA} por cada participante, actualizando su saldo (\mbox{RN-08}). \textit{(Ver \mbox{EX-01})}
      \item El Servidor de partidas cierra las \textit{Participacion}, dejándolas sin terminar en ningún caso, de modo que sus jugadores vuelvan a poder buscar partida (\mbox{RN-09}).
      \item El Servidor de partidas confirma la transacción. \textit{(Ver \mbox{EX-02})}
      \item El Servidor de partidas notifica a cada cliente con \texttt{MATCH\_\allowbreak{}ENDED}, incluyendo el resultado, la forma de término, el cambio de elo, la experiencia y las monedas obtenidas.
      \item El SJM despliega la \texttt{GUI\_\allowbreak{}MatchEnd} con el resumen sobre el tablero, que sigue visible detrás.
    \end{cuenum}} \\
   & \cuCelda{\begin{cuenum}[start=17]
      \item Si hubo subida de nivel o cambio de división, el SJM despliega sus pantallas propias con lo desbloqueado. \textit{(Ver \mbox{FA-02}, \mbox{FA-03})}
      \item El SJM ofrece las opciones ``Buscar otra'', ``Ver repetición'', ``Añadir al rival'' y ``Volver al menú''.
      \item Fin del caso de uso.
    \end{cuenum}} \\ \hline
  \textbf{Flujos alternos} & \cuCelda{\textbf{\mbox{FA-01}: La partida no es clasificatoria}\par
    \begin{cuenum}[start=1]
      \item El Servidor de partidas omite los pasos 7, 8, 10 y 12 del FN: no hay elo, ni división, ni monedas (\mbox{RN-03}).
      \item El Servidor de partidas actualiza de todos modos la \textit{EstadisticaModo} solo en lo que no depende del elo, si el \textit{Modo} lo contempla, y otorga la experiencia del paso 11.
      \item Si la partida era contra la IA, el Servidor de partidas prepara además el repaso de tres puntos con los errores del jugador, calculado sobre las \textit{Jugada} registradas.
      \item El flujo continúa en el paso 13 del FN.
    \end{cuenum}} \\
   & \cuCelda{\textbf{\mbox{FA-02}: Un participante cambia de división}\par
    \begin{cuenum}[start=1]
      \item El Servidor de partidas escribe la división nueva en la \textit{EstadisticaModo} y crea un \textit{HistorialDivision} con la división anterior, la nueva y la fecha, para que la gráfica de evolución del \mbox{CU-36} pueda marcarlo.
      \item El SJM despliega la pantalla de cambio de división con las recompensas desbloqueadas, si las hay.
      \item El flujo continúa en el paso 18 del FN.
    \end{cuenum}} \\
   & \cuCelda{\textbf{\mbox{FA-03}: Un participante sube de nivel}\par
    \begin{cuenum}[start=1]
      \item El Servidor de partidas escribe el nivel nuevo en el \textit{Usuario} y otorga las recompensas asociadas, creando los \textit{UsuarioObjeto} correspondientes, una \textit{Caja} de recompensa con origen \texttt{PARTIDA} y caducidad a catorce días cuando la recompensa lo indica, y, si procede, un \textit{MovimientoMoneda} de tipo \texttt{NIVEL}.
      \item El SJM despliega la pantalla de subida de nivel con lo desbloqueado.
      \item El flujo continúa en el paso 18 del FN.
    \end{cuenum}} \\
   & \cuCelda{\textbf{\mbox{FA-04}: La partida termina por abandono en el modo de cuatro jugadores}\par
    \begin{cuenum}[start=1]
      \item El Servidor de partidas marca la \textit{Participacion} del que abandonó con resultado perdido y forma de término \texttt{ABANDONO}, y retira su peón del tablero.
      \item Los muros que ese jugador había colocado permanecen en el tablero (\mbox{RN-10}).
      \item Si quedan al menos dos jugadores activos, la partida \textbf{no} termina: el flujo vuelve al \mbox{CU-20} y este caso solo se aplicó a esa participación. En caso contrario, el flujo continúa en el paso 1 del FN.
    \end{cuenum}} \\
   & \cuCelda{\textbf{\mbox{FA-05}: La partida termina en tablas}\par
    \begin{cuenum}[start=1]
      \item El Servidor de partidas escribe el resultado de tablas en todas las \textit{Participacion}.
      \item El Servidor de partidas calcula el elo conforme a \mbox{RN-04} tratando el resultado como intermedio, de modo que el de menor elo gane algo y el de mayor elo pierda algo.
      \item La \textit{EstadisticaModo} incrementa \texttt{partidas\_\allowbreak{}jugadas} sin incrementar ganadas ni perdidas, y la \texttt{racha\_\allowbreak{}actual} se rompe sin contar como derrota (\mbox{RN-06}).
      \item El flujo continúa en el paso 10 del FN.
    \end{cuenum}} \\
   & \cuCelda{\textbf{\mbox{FA-06}: La partida se gana por reloj}\par
    \begin{cuenum}[start=1]
      \item El Servidor de partidas registra la forma de término \texttt{TIEMPO\_\allowbreak{}AGOTADO}.
      \item El Servidor de partidas no cuenta esta victoria para la racha limpia, que solo admite victorias por llegada a meta (\mbox{RN-11}).
      \item El flujo continúa en el paso 7 del FN.
    \end{cuenum}} \\
   & \cuCelda{\textbf{\mbox{FA-07}: Un participante tiene la sesión cerrada al terminar la partida}\par
    \begin{cuenum}[start=1]
      \item El Servidor de partidas aplica igualmente todos los pasos del FN sobre ese participante: la progresión no depende de que esté conectado.
      \item El Servidor de partidas no le envía \texttt{MATCH\_\allowbreak{}ENDED}.
      \item El jugador verá el resultado en su historial la próxima vez que entre.
      \item El flujo continúa en el paso 16 del FN para los demás.
    \end{cuenum}} \\
   & \cuCelda{\textbf{\mbox{FA-08}: La partida tenía enlaces de espectador o venía de una sala}\par
    \begin{cuenum}[start=1]
      \item El Servidor de partidas marca con \texttt{activo} en falso los \textit{EnlaceEspectador} de la partida, dentro de la transacción del cierre (\mbox{RN-06} del \mbox{CU-34}).
      \item Si la partida nació de una \textit{Sala}, el Servidor de partidas cambia su \texttt{estado} a \texttt{CERRADA} y elimina sus \textit{SalaParticipante}.
      \item El Servidor de partidas retira de la difusión a los espectadores después de enviarles el resultado con el retraso previsto.
      \item El flujo continúa en el paso 14 del FN.
    \end{cuenum}} \\ \hline
  \textbf{Excepciones} & \cuCelda{\textbf{\mbox{EX-01}: Fallo al escribir en la base de datos}\par
    \begin{cuenum}[start=1]
      \item El Servidor de partidas revierte la transacción completa, de modo que no quede una partida cerrada con estadísticas a medio actualizar, ni monedas otorgadas sin resultado escrito (\mbox{RN-12}).
      \item El Servidor de partidas registra la excepción en la bitácora del servidor con un identificador de incidencia.
      \item El Servidor de partidas deja la \textit{Partida} en \texttt{EN\_\allowbreak{}CURSO} y reintenta el cierre conforme a su política de reintentos.
      \item Si los reintentos se agotan, el Servidor de partidas marca la \textit{Partida} como \texttt{PENDIENTE\_\allowbreak{}DE\_\allowbreak{}CIERRE} y notifica a los clientes que el resultado se aplicará más tarde (\mbox{RN-13}).
      \item Fin del caso de uso.
    \end{cuenum}} \\
   & \cuCelda{\textbf{\mbox{EX-02}: Fallo al confirmar la transacción}\par
    \begin{cuenum}[start=1]
      \item El Servidor de partidas trata la transacción como fallida y no notifica el fin a nadie.
      \item El Servidor de partidas registra la excepción señalando que el estado del cierre es indeterminado.
      \item El flujo continúa en el paso 3 del \mbox{EX-01}.
    \end{cuenum}} \\
   & \cuCelda{\textbf{\mbox{EX-03}: Un cliente no recibe la notificación de fin}\par
    \begin{cuenum}[start=1]
      \item El Servidor de partidas no reintenta indefinidamente: el resultado ya está escrito y es consultable.
      \item Al reconectar, el SJM detecta que su \textit{Participacion} terminó y despliega la \texttt{GUI\_\allowbreak{}MatchEnd} con el resultado leído del servidor.
      \item Fin del caso de uso.
    \end{cuenum}} \\
   & \cuCelda{\textbf{\mbox{EX-04}: El cálculo de elo no puede completarse}\par
    \begin{cuenum}[start=1]
      \item El Servidor de partidas detecta que falta el elo inicial en alguna \textit{Participacion}, lo que solo puede ocurrir por una partida creada de forma incorrecta.
      \item El Servidor de partidas cierra la partida sin aplicar elo ni monedas, y registra la incidencia en la bitácora del servidor para su revisión.
      \item El flujo continúa en el paso 13 del FN.
    \end{cuenum}} \\ \hline
  \textbf{Postcondiciones} & \cuCelda{\begin{cureglas}
      \item \textbf{\mbox{POST-01}} La \textit{Partida} tiene estado \texttt{FINALIZADA}, su \texttt{fecha\_\allowbreak{}fin}, su forma de término y su ganador, si lo hay.
      \item \textbf{\mbox{POST-02}} Cada \textit{Participacion} tiene su resultado, su posición, su elo inicial, su elo final y la diferencia entre ambos.
      \item \textbf{\mbox{POST-03}} Ninguna \textit{Participacion} de la partida sigue sin terminar, de modo que todos sus jugadores pueden volver a buscar partida (\mbox{D-01}).
      \item \textbf{\mbox{POST-04}} Si la partida era clasificatoria, la \textit{EstadisticaModo} de cada participante en ese \textit{Modo} refleja la partida en todos sus contadores.
      \item \textbf{\mbox{POST-05}} Si la partida era clasificatoria, cada participante tiene un \textit{MovimientoMoneda} de tipo \texttt{PARTIDA} y su saldo actualizado.
      \item \textbf{\mbox{POST-06}} Si algún participante cruzó un umbral, su división quedó escrita y el cambio registrado con su fecha.
      \item \textbf{\mbox{POST-07}} Si la partida no era clasificatoria, ningún elo, ninguna división y ningún saldo de monedas cambió.
    \end{cureglas}} \\
   & \cuCelda{\begin{cureglas}
      \item \textbf{\mbox{POST-08}} Si el caso termina por una excepción, la partida no quedó cerrada a medias: o se aplicó todo, o sigue \texttt{EN\_\allowbreak{}CURSO} o \texttt{PENDIENTE\_\allowbreak{}DE\_\allowbreak{}CIERRE}.
      \item \textbf{\mbox{POST-09}} Las \textit{Jugada} de la partida siguen íntegras y permiten la repetición del \mbox{CU-37}.
    \end{cureglas}} \\ \hline
  \textbf{Reglas de negocio} & \cuCelda{\begin{cureglas}
      \item \textbf{\mbox{RN-01}} Una partida termina de cinco formas: llegada a meta, reloj agotado, rendición, tablas aceptadas y abandono. La forma de término se guarda, no solo el resultado, porque no todas valen lo mismo.
      \item \textbf{\mbox{RN-02}} En el modo de cuatro jugadores, quien llega a meta queda clasificado con su posición y la partida sigue para los demás.
      \item \textbf{\mbox{RN-03}} Solo las partidas clasificatorias otorgan elo y monedas. Las de sala privada y las de la IA otorgan experiencia pero no alteran el ranking ni la economía.
      \item \textbf{\mbox{RN-04}} El elo se calcula a partir del elo de entrada de cada participante, no del que tenga en ese momento en su cuenta. Dos partidas simultáneas no pueden ocurrir (\mbox{D-01}), pero el elo de entrada es de todos modos lo único que hace reproducible el cálculo.
      \item \textbf{\mbox{RN-05}} El cambio de elo se guarda en la \textit{Participacion}. El historial del \mbox{CU-36} lo muestra partida por partida y no podría recomponerlo restando estados sucesivos de la cuenta.
      \item \textbf{\mbox{RN-06}} Los contadores de \textit{EstadisticaModo} son por modo, no globales. El elo del menú y del ranking general es el del modo clásico 9×9.
      \item \textbf{\mbox{RN-07}} Los cambios de división respetan un umbral de descenso visible y unas partidas de margen, de modo que una sola derrota no haga bajar de división.
    \end{cureglas}} \\
   & \cuCelda{\begin{cureglas}
      \item \textbf{\mbox{RN-08}} Toda entrada o salida de monedas deja un \textit{MovimientoMoneda}. El saldo del \textit{Usuario} es la suma de sus movimientos y no una cifra independiente de ellos.
      \item \textbf{\mbox{RN-09}} Cerrar la partida cierra todas sus participaciones. Mientras una siga sin terminar, su jugador no puede entrar a otra partida.
      \item \textbf{\mbox{RN-10}} Los muros de un jugador que abandonó permanecen en el tablero. El abandono retira al jugador, no lo que ya hizo.
      \item \textbf{\mbox{RN-11}} Una victoria por reloj se distingue de una victoria por llegada a meta y no cuenta para la racha limpia.
      \item \textbf{\mbox{RN-12}} El cierre de la partida, la escritura de resultados, la actualización de estadísticas y los movimientos de monedas ocurren en una sola transacción.
      \item \textbf{\mbox{RN-13}} Una partida que no puede cerrarse queda marcada como pendiente de cierre y se reintenta. No se descarta: los jugadores quedarían bloqueados para siempre por el \mbox{RN-09}.
      \item \textbf{\mbox{RN-14}} Hacen falta cinco partidas en un modo para aparecer en su clasificación.
    \end{cureglas}} \\ \hline
  \textbf{Observación} & \cuCelda{\textit{Sobre por qué este caso es el único que escribe el elo.}\par
    Es el caso de uso más caro del sistema y el que más tablas toca de una vez, y por eso conviene decir en voz alta por qué no se reparte. El elo, las monedas y los contadores dependen todos del mismo resultado y deben moverse juntos o no moverse: un jugador con el elo actualizado y las monedas sin otorgar es un error que nadie sabría reparar después. Repartirlo entre varios casos —uno que cierre, otro que puntúe, otro que pague— obligaría a coordinarlos con una transacción que los abarcara a todos, que es exactamente lo que este caso ya es.} \\ \hline
  \textbf{Datos que toca} & \cuCelda{\begin{cudatos}
      \item \textbf{\textit{Partida}:} escribe estado, \texttt{fecha\_\allowbreak{}fin}, forma de término, ganador \textit{(paso 5)}
      \item \textbf{\textit{Participacion}:} escribe resultado, posición, elo final, diferencia, reloj restante \textit{(pasos 6, 8, 13)}
      \item \textbf{\textit{Jugada}:} lee, para contar muros colocados y para el repaso de la IA \textit{(pasos 9, \mbox{FA-01})}
      \item \textbf{\textit{EstadisticaModo}:} escribe todos los contadores y la división \textit{(pasos 9, \mbox{FA-02})}
      \item \textbf{\textit{HistorialDivision}:} crea uno por cambio de división \textit{(paso FA-02)}
      \item \textbf{\textit{Usuario}:} escribe experiencia, nivel y saldo de monedas \textit{(pasos 11, 12, \mbox{FA-03})}
      \item \textbf{\textit{MovimientoMoneda}:} crea uno por participante \textit{(pasos 12, \mbox{FA-03})}
      \item \textbf{\textit{UsuarioObjeto}:} crea, por las recompensas de nivel \textit{(paso \mbox{FA-03})}
      \item \textbf{\textit{Caja}:} crea las de recompensa \textit{(paso FA-03)}
      \item \textbf{\textit{EnlaceEspectador}:} desactiva los de la partida \textit{(paso FA-08)}
      \item \textbf{\textit{Sala}, \textit{SalaParticipante}:} cierra la sala de origen y elimina sus miembros \textit{(paso FA-08)}
      \item \textbf{\textit{Modo}:} lee, para los parámetros de puntuación \textit{(paso 7)}
    \end{cudatos}} \\ \hline
\end{fichacu}
\clearpage
